\section{Exposure Adjustment and Spatial Closure}\label{sec:spatial-closure}

How does peer adjustment transform the stand-alone exposures $\xi$ into
peer-adjusted exposures $B$?
We answer with an as-if reduced-form adjustment criterion that summarizes costly
operational, financing, or portfolio reconfiguration.
The criterion balances fidelity to the firm's stand-alone exposure against
alignment with its peer-weighted exposure.
We call the resulting link \emph{spatial closure}: the objective yields, rather
than assumes, a spatial autoregression in exposures.

To formalize this trade-off, fix an interaction matrix $W$ satisfying
\Cref{def:interaction-matrix} and an adjustment intensity $\lambda\geq0$.
For a fixed peer profile $b_{-i}$, let $a\in\mathcal H$ denote firm $i$'s
candidate exposure, let $\xi_i$ denote its stand-alone exposure, and let
$b_j\in\mathcal H$ denote the peer exposures.
The adjustment criterion is the following quadratic problem.

\begin{definition}[Quadratic adjustment problem]\label{def:adjustment-problem}
  Holding the profile $b_{-i}$ of peer exposures fixed, firm $i$ chooses
  $a\in\mathcal H$ to minimize
  \begin{equation}\label{eq:adjustment-objective}
    \frac{1}{2}\norm{a-\xi_i}^2
    +\frac{\lambda}{2}\sum_j W_{ij}\norm{a-b_j}^2 .
  \end{equation}
  A profile $b$ is an equilibrium when every $b_i$ minimizes
  \eqref{eq:adjustment-objective} at $b_{-i}$.
\end{definition}

Both terms are measured in squared exposure units.
The first penalizes departure from the stand-alone exposure, whereas the second
penalizes disagreement with the peer profile.
Row stochasticity keeps the scale of the peer penalty comparable across firms,
and $\lambda$ measures its weight relative to the stand-alone-exposure penalty.
The quadratic form and fixed $W$ are maintained inputs to the model.

To characterize equilibrium, first ask when each firm's candidate exposure
minimizes its criterion.
Unit row sums collect the quadratic terms and yield a condition that is both
necessary and sufficient.

\begin{lemma}[Adjustment equilibrium
  condition]\label{thm:adjustment-equilibrium}
  For $\lambda\geq0$ and row-stochastic $W$, a profile $b$ is an equilibrium of
  \Cref{def:adjustment-problem} if and only if
  \begin{equation}\label{eq:adjustment-stationarity}
    b_i-\xi_i+\lambda\sum_j W_{ij}(b_i-b_j)=0
    \qquad\text{for every }i.
  \end{equation}
\end{lemma}

Economically, \eqref{eq:adjustment-stationarity} balances displacement from the
stand-alone exposure against the weighted gap from peer exposures.
Mathematically, completing the square isolates a linear term whose coefficient
is the left-hand side of \eqref{eq:adjustment-stationarity}.
A minimizer forces that coefficient to vanish; once it does, the objective gap
at any alternative $a$ is
$\tfrac{1}{2}(1+\lambda)\norm{a-b_i}^2$.
The condition is therefore necessary and sufficient in any real inner-product
space, without a finite-dimensional differentiability argument.
This lemma characterizes equilibrium for a fixed admissible $W$.

To close the firm-level conditions as a simultaneous system, use
$\sum_{j} W_{ij}=1$ in \eqref{eq:adjustment-stationarity} and rearrange.

\begin{theorem}[Spatial closure]\label{thm:spatial-closure}
  Under the hypotheses of \Cref{thm:adjustment-equilibrium}, every equilibrium
  satisfies the Hilbert-valued spatial autoregression
  \begin{equation}\label{eq:hilbert-sar}
    B=\rho\,WB+(1-\rho)\,\xi,
    \qquad
    \rho=\frac{\lambda}{1+\lambda}.
  \end{equation}
\end{theorem}

Equation~\eqref{eq:hilbert-sar} is the spatial closure: the familiar spatial lag
follows from the two adjustment penalties rather than entering as an assumed
return equation.
For each firm, peer-adjusted exposure is a convex balance between its peer
average and its stand-alone exposure.
Accordingly, $\rho$ is a model-implied index of relative adjustment, not an
identified causal share attributable to peers.
The map $\lambda\mapsto\lambda/(1+\lambda)$ is a strictly increasing bijection
from nonnegative adjustment intensity to $0\leq\rho<1$.
Consequently, the relative penalty weight can be recovered from a spatial
coefficient by

\begin{equation}\label{eq:lambda-from-rho}
  \lambda=\frac{\rho}{1-\rho}.
\end{equation}

A value $\rho=1/2$ gives the peer average and the stand-alone exposure equal
weight, while larger values place more weight on peers.
At $\rho=0$, peer-adjusted exposure equals stand-alone exposure, whereas values
approaching one place progressively greater model-implied weight on peer
alignment.

To express the simultaneous system as a unique reduced form, the spatial
multiplier must exist.
We impose the following standard sufficient stability condition.

\begin{hypothesis}[Stability condition]\label{hyp:spectral-gate}
  The pair $(\rho, W)$ satisfies $\norm{\rho W} < 1$ in the
  $\ell^\infty$-induced operator norm.
  That norm is the largest absolute row sum, so the nonnegative unit rows of
  \Cref{def:interaction-matrix} give $\norm{W}=1$ and $0\leq\rho<1$ is
  sufficient.
\end{hypothesis}

Economically, the condition keeps iterated peer feedback anchored by the
stand-alone exposures.
Mathematically, rearranging \eqref{eq:hilbert-sar} and applying the convergent
Neumann series for $(I-\rho W)^{-1}$ gives the following unique equilibrium.

\begin{theorem}[Reduced form]\label{thm:sar-reduced-form}
  If in addition $\norm{\rho W}<1$, the equilibrium is unique and
  \begin{equation}\label{eq:hilbert-reduced-form}
    B=(1-\rho)(I-\rho W)^{-1}\xi .
  \end{equation}
\end{theorem}

The multiplier aggregates direct and iterated peer adjustment, while the factor
$1-\rho$ preserves the scale of the stand-alone exposures.
This is an equilibrium result for exposures; its observed-return interpretation
requires the projection argument in \Cref{sec:identification-estimation}.

To connect spatial closure to the pairwise exposure model of
\citet{gawronsky_continuous_2026}, consider the zero-feedback benchmark.

\begin{corollary}[Zero-feedback nesting]\label{cor:zero-feedback-nesting}
  At $\rho=0$, $B=\xi$.
\end{corollary}

The equality is exact rather than limiting: when peer alignment receives zero
weight, peer-adjusted exposure equals stand-alone exposure.
The benchmark therefore recovers the exposure object used by their pairwise
covariance restriction without reproducing that theory.

\subsection{Joint Adjustment Across Two Interaction Fields}
\label{sec:two-field-closure}

Nothing in the adjustment criterion requires the researcher to settle on one
definition of a peer.
Distributional similarity and explicit news links may each carry distinct
conditional information about peer relevance.
The economic question is therefore how peer-adjusted exposure balances both
peer averages against the firm's stand-alone exposure.
The two-field objective represents that trade-off directly.

Let $W^{B}$ and $W^{N}$ be two matrices satisfying
\Cref{def:interaction-matrix}.
Let $\lambda_B,\lambda_N\geq0$ be the relative penalty weights on misalignment
with each.
The labels anticipate the two fields the application supplies; the statements
below use nothing beyond admissibility of each matrix.

\begin{definition}[Two-field adjustment problem]\label{def:two-field-adjustment}
  Holding the profile $b_{-i}$ fixed, firm $i$ chooses $a\in\mathcal H$ to
  minimize
  \begin{equation}\label{eq:two-field-objective}
    \frac{1}{2}\norm{a-\xi_i}^2
    +\frac{\lambda_B}{2}\sum_j W^{B}_{ij}\norm{a-b_j}^2
    +\frac{\lambda_N}{2}\sum_j W^{N}_{ij}\norm{a-b_j}^2 .
  \end{equation}
  A profile $b$ is an equilibrium when every $b_i$ minimizes
  \eqref{eq:two-field-objective} at $b_{-i}$.
\end{definition}

All three terms carry squared exposure units, so $\lambda_B$ and $\lambda_N$
are dimensionless penalty weights relative to the stand-alone term.
Setting either intensity to zero returns \Cref{def:adjustment-problem} at the
other field.

To connect the two-field objective to the one-field closure, write
$\lambda=\lambda_B+\lambda_N$ and, when $\lambda>0$,
$\theta=\lambda_B/\lambda$, and define the mixture
$W(\theta)=\theta W^{B}+(1-\theta)W^{N}$.
The two peer penalties in \eqref{eq:two-field-objective} then equal the
one-field penalty in \eqref{eq:adjustment-objective} at total intensity
$\lambda$ and matrix $W(\theta)$.
Before applying the one-field results, the next result verifies that this
mixture remains an admissible peer matrix.

\begin{corollary}[Mixture nesting]\label{cor:mixture-nesting}
  For $0\leq\theta\leq1$ and $W^{B},W^{N}$ satisfying
  \Cref{def:interaction-matrix}, the mixture $W(\theta)$ satisfies
  \Cref{def:interaction-matrix}.
  The endpoints $\theta=1$ and $\theta=0$
  return $W^{B}$ and $W^{N}$.
\end{corollary}

Economically, the mixture allocates one unit of peer weight across two peer
definitions while preserving the peer-average interpretation.
Nonnegativity and the zero diagonal are preserved entrywise, and the row sums
are a convex combination of two unit row sums.
Thus the one-field results apply unchanged.

With admissibility established, the joint-field result answers how total
model-implied adjustment divides across the two channels.

\begin{theorem}[Two-field spatial closure]\label{thm:two-field-closure}
  For $\lambda_B,\lambda_N\geq0$ and admissible $W^{B},W^{N}$, every
  equilibrium of \Cref{def:two-field-adjustment} satisfies
  \begin{equation}\label{eq:two-field-sar}
    B=\rho_B\,W^{B}B+\rho_N\,W^{N}B+(1-\rho_B-\rho_N)\,\xi,
    \qquad
    \rho_k=\frac{\lambda_k}{1+\lambda_B+\lambda_N},
  \end{equation}
  for $k\in\{B,N\}$.
\end{theorem}

Equation~\eqref{eq:two-field-sar} makes the channel allocation explicit:
$\rho_B+\rho_N$ is the total model-implied peer-adjustment weight, while
$\rho_B$ and $\rho_N$ assign that total to the two fields.
The remaining weight $1-\rho_B-\rho_N$ anchors peer-adjusted exposure to the
firm's stand-alone exposure.
When $\lambda_N>0$, equivalently $\rho_N>0$, the ratio
$\lambda_B/\lambda_N=\rho_B/\rho_N$ describes how peer adjustment divides
between the channels.
To obtain these coefficients, apply \Cref{thm:spatial-closure} at
$(\lambda,W(\theta))$ and expand
$B=\rho\,W(\theta)B+(1-\rho)\xi$ with $\rho=\lambda/(1+\lambda)$.

To establish stability of the joint system and recover each channel's relative
penalty weight, combine the coefficient levels in the following result.

\begin{proposition}[Two-field stability and channel inversion]
  \label{prop:channel-inversion}
  Under the hypotheses of \Cref{thm:two-field-closure},
  $\rho_B+\rho_N=\lambda/(1+\lambda)<1$, and
  $\norm{\rho_B\,W^{B}+\rho_N\,W^{N}}\leq\rho_B+\rho_N$ in the
  $\ell^\infty$-induced operator norm, so \Cref{hyp:spectral-gate} holds and
  the reduced form of \Cref{thm:sar-reduced-form} applies at the mixture.
  Each intensity is recovered from the coefficients by
  \begin{equation}\label{eq:channel-inversion}
    \lambda_k=\frac{\rho_k}{1-\rho_B-\rho_N},
    \qquad k\in\{B,N\}.
  \end{equation}
\end{proposition}

Equation~\eqref{eq:channel-inversion} generalizes
\eqref{eq:lambda-from-rho}: each channel's adjustment index is its own spatial
coefficient measured against the weight the firm still places on its
stand-alone exposure.
The norm bound follows from \Cref{cor:mixture-nesting} and
$\norm{W(\theta)}\leq1$.
For the inversion, $1-\rho_B-\rho_N=(1+\lambda)^{-1}$, so dividing $\rho_k$ by
this common stand-alone weight returns $\lambda_k$.

The closure argument is complete conditional on an admissible interaction
field.
Where does that field come from when peer relevance must be inferred from
distributions of firm text?
The next section develops a \emph{target-anchored Wasserstein barycentric
reconstruction}: candidate firms are aligned to a fixed target, and their
distributional spanning weights form the \emph{barycentric interaction field}
$W^\flat$.
Each row is admissible, and the computation does not require a conventional
kernel bandwidth.
